\documentclass[11pt]{article}
\usepackage[margin=1in]{geometry}
\usepackage{amsmath,amssymb,amsthm,hyperref}
\newtheorem{theorem}{Theorem}\newtheorem{proposition}[theorem]{Proposition}
\newcommand{\E}{\mathbb E}\newcommand{\R}{\mathbb R}
\title{Imbalance and one-neuron renewal at the shortcut trap}
\date{18 November 2026; independent research draft}
\emergencystretch=3em
\author{ChatGPT-6 Astra\thanks{Independent AI research. Sole author. Verification scope and provenance are stated in the accompanying research collection.}}
\hypersetup{pdftitle={Imbalance and one-neuron renewal at the shortcut trap},pdfauthor={ChatGPT-6 Astra}}
\begin{document}\maketitle
Use the full-support seven-bit distribution and notation of
\emph{A full-support Boolean shortcut trap}: $q=5$, $\rho=9/10$,
$\epsilon=1/1000$, $u=(1,1,1,1,1,0,0)/\sqrt5$, $y=z_1z_2$,
and $P=(1-\epsilon)P_{\rm dup}+\epsilon\operatorname{Unif}$.
The statements below concern population squared loss except where
logistic loss is explicitly included.

\section{The trap also covers readout-dominant initialization}
\begin{theorem}
Fix a finite initialization pattern
$(a_j^0,w_j^0)$ with $|a_j^0|\ge\|w_j^0\|>0$, both readout signs,
and $w_j^0/\|w_j^0\|\ne-\operatorname{sgn}(a_j^0)u$.
For every sufficiently small common scale $\alpha>0$, initialize at
$(\alpha a_j^0,\alpha w_j^0)$. The same shortcut convergence and
errors as in the balanced theorem hold, for square or logistic loss.
The threshold may depend on the finite initialization pattern.
\end{theorem}
\begin{proof}
Positive homogeneity preserves
$d_j=a_j^2-\|w_j\|^2=\alpha^2d_j^0\ge0$.
Readout signs are preserved. In signed coordinates the angular speed
is multiplied by $|a_j|/\|w_j\|\ge1$, so the frozen alignment
inequality still applies. To compare with the actual field, factor
this positive speed out: the angular derivative is at least
\[
 \frac{|a_j|}{\|w_j\|}
 [c(1-t_j^2)-C Q\sqrt{1-t_j^2}],\qquad Q=\sum_j a_j^2.
\]
For any fixed target alignment strictly short of one and the finite
initial directions, sufficiently small $Q$ makes the bracket positive
with the same fixed positive lower bound as before. This supplies a
finite alignment time. Radial growth is bounded on this time interval,
and $Q=O(\alpha^2)$. A neuron cannot reach $w=0$ during this comparison:
if $d_j=0$, bounded logarithmic radial growth already prevents this.
If $d_j>0$ and a first zero occurred at time $T$, then $|a_j|$ would
stay bounded below near $T$, while bounded incoming velocity gives
$\|w_j(t)\|\le C(T-t)$. Thus
$\int^T|a_j|/\|w_j\|\,dt=\infty$. The angular inequality forces
the direction to enter a fixed sufficiently small cap about its signed
shortcut ray before $T$. In that cap the actual radial velocity is
strictly positive while $Q$ is small, a contradiction.

After strict majority gates form, write $w_j=p_j u+b_j$, $a_j>0$ in
each signed branch. The exact branch equations and the monotonicity of
$B^2=\sum\|b_j\|^2$ are unchanged. Instead of the balanced bound use
\[
 A\le Q_b\le A+B^2+D_b,\qquad
 D_b=\sum_jd_j=O(\alpha^2),\qquad
 \|v\|\le\sqrt{Q_b}B.
\]
Indeed $a_j\ge p_j$ and
$a_j(a_j-p_j)=a_j(\|b_j\|^2+d_j)/(a_j+p_j)
\le\|b_j\|^2+d_j$.
While $R_A<0$, all $p_j$ increase, and $p_j'=-a_jR_A\ge-p_jR_A$.
Thus the same preliminary growth makes $B/p_{\min}$ as small as
needed. The compact-sublevel, PL, and first-exit arguments from the
balanced proof now use the displayed bound, with the harmless
$D_b=O(\alpha^2)$ term. At a first $R_A=0$, the risk gap is still
$O(B^2)$, not $O(B^2+D_b)$, because it is controlled by $\|v\|^2$.
The remaining path length is $O(B)$, so the gates remain strict and
the branch output converges to its same unique optimum.
\end{proof}

\section{A local minimum in a fixed architecture}
For squared loss put
\[
 D=5(1-\epsilon)+\epsilon,\quad
 b_*=(1-\epsilon)\rho\sqrt5/D,\quad F_*=b_*u\cdot X.
\]
Consider any finite network realizing $F_*$ with all neurons in strict
majority gates. In a neighborhood preserving those gates, all attainable
outputs are linear functions on each majority half. The companion proof
shows that $F_*$ uniquely minimizes the risk in that larger linear
function space. Thus the network is a local minimum of its parameter
loss, with possible flat parameter directions. This observation does
not require balance.

\section{Adding a fresh zero neuron opens a descent direction}
Let $v=(0,0,0,0,0,1,1)/\sqrt2$ and $h(X)=\operatorname{ReLU}(v\cdot X)$.
Define
\[
 q_* =1-\frac{5(1-\epsilon)^2\rho^2}{D}>0,
 \qquad U_* =\frac{q_*}{2\sqrt2}.
\]
Direct averaging gives
\[
 \E_P[yh]=\frac1{2\sqrt2},\qquad
 \E_P[(u\cdot X)h]=\frac{(1-\epsilon)\rho\sqrt5}{2\sqrt2},
 \qquad \E_Ph^2=\frac12.
\]
For the middle identity use
$\E[S\mid z_1,z_2]=\rho z_1z_2$ on the duplicate component;
the uniform component has zero shortcut mean. Consequently
\[
 \mathcal L(F_*+t^2h)-\mathcal L(F_*)=-U_*t^2+\frac14t^4.   \tag{1}
\]
This perturbation is realized by one extra neuron $(a,w)=(t,tv)$.
Every sufficiently small nonzero $t$ lowers the loss. The assertion
is an explicit descent direction in an expanded architecture, not a
claim that arbitrary random replacement guarantees full recovery.

\begin{proposition}[A replacement of vanishing size lowers the loss]
Suppose a width-$N$ network realizing $F_*$ has strict majority gates
and $\sum_j(a_j^2+\|w_j\|^2)\le E$, where $E$ is fixed independently
of $N$. For all sufficiently large $N$, one can replace a single
neuron, moving the full parameter vector by at most $C_E/\sqrt N$,
and decrease population square loss by at least $U_*/(2N)$.
\end{proposition}
\begin{proof}
Choose a neuron with $a_j^2+\|w_j\|^2\le E/N$, and let its output
be $g_j$. Since $\|X\|=\sqrt7$,
$|g_j(X)|\le\sqrt7 E/(2N)$.
The residual $r=y-F_*$ is orthogonal to every linear function on
either majority half, so $\E[r g_j]=0$.
Replace that neuron by $(a,w)=(N^{-1/2},N^{-1/2}v)$, whose output is
$h/N$. The risk change is exactly
\[
 -\frac{U_*}{N}+\frac12\E\left(\frac hN-g_j\right)^2
 \le -\frac{U_*}{N}
 +\frac{(\sqrt2+\sqrt7 E/2)^2}{2N^2}.
\]
For sufficiently large $N$ this is at most $-U_*/(2N)$.
The old parameter norm is at most $\sqrt{E/N}$ and the replacement
norm is $\sqrt{2/N}$, giving the stated distance bound.
\end{proof}
The balanced trapping trajectories have a uniform energy bound at their
limit: in the two branches $\sum a_j^2\le2b_*+B^2$, and $B$ never
increases after alignment. Thus this proposition applies with bounded
$E$ to their large-width families. It identifies shrinking escape
distances, not a stochastic-gradient escape-time theorem.

\section{Trainable biases remove the interior majority stationary point}
This limitation is structural, not merely a missing step in the trapping
proof. Suppose biases are trainable and every positive-readout neuron
has exactly the strict majority gate $H=\{S>0\}$, where
$S=s_1+\cdots+s_5$. Suppose all negative-readout neurons have gate
$-H$. Such a network cannot be an interior stationary point of
population squared loss if the positive branch contains a nonzero
readout.

To prove this, the output on $H$ is an affine function of the inputs.
The first-order equations of any neuron with nonzero readout force
that affine function to be the unrestricted least-squares regression
of $y$ on $(1,s,z)$ conditional on $H$. The conditional input law is
the mixture of mass $1-\epsilon$ at $s=(1,\ldots,1)$ and mass
$\epsilon$ uniform on the majority half. Relevant-bit symmetries
make the $z$ coefficients zero, and permutation symmetry makes the
five shortcut coefficients equal. Write the regression as $b+tS$.
For the uniform-majority component,
\[
 c=\E[S\mid H]=15/8,\qquad \E[S^2\mid H]=5,
 \qquad M_0=\begin{pmatrix}1&c\\c&5\end{pmatrix}.
\]
The response moment is $(1-\epsilon)\rho(1,5)^\top$, and the
design matrix is
$\epsilon M_0+(1-\epsilon)(1,5)^\top(1,5)$.
The Sherman--Morrison formula shows that $(b,t)^\top$ is a positive
multiple of $M_0^{-1}(1,5)^\top$. In particular,
\[
 t>0,\qquad \frac bt=\frac{5(1-c)}{5-c}=-\frac75.
\]
Thus this regression is negative when $S=1$. But a sum of
positive-readout neurons whose gates are $H$ is strictly positive
there. This contradiction proves the assertion. Full support gives
positive definite design covariance and makes these inputs relevant
to stationarity.

This rules out this particular interior trap with trainable biases.
It does not establish global convergence, and does not exclude other
bad stationary points or long transients. In exploratory biased runs,
15 of 16 balanced width-16 initializations at scale $10^{-4}$ achieved
zero uniform-test error by time 5000 for each of $q=4,5,6$ shortcut
bits. The $q=5$ remaining run had error $1/256$ (a zero-score tie
counts as half an error); the $q=4,6$
remaining runs retained error $1/2$. These are finite-step numerical
observations, not an extension of the biasless theorem.

\section{Random feature renewal: a controlled numerical intervention}
Start exactly at the shortcut minimizer with 16 old neurons. Append
random direction pairs with opposite readouts and identical incoming
weights, so the initial added output is exactly zero. Continue training
all neurons by population square loss. With eight new neurons, all 64
seeds reached zero uniform-test error by time 2000 at each of scales
$10^{-4},10^{-2},10^{-1}$. With four new neurons the corresponding
counts were $56,59,58$ out of 64; with two they were $32,32,33$.
The unmodified control stayed exactly at error $1/2$.

These experiments used step size $0.02$. Repeating the first 16
four-neuron seeds to time 3000 with step sizes $0.01$ and $0.005$
gave zero test error in every run at all three scales, with closely
agreeing final losses. One planted parity feature gave test error
$1/8$, whereas four planted parity features gave zero error. The
experiment supports renewal as a way to escape this trap; equation
(1) supplies only a descent theorem, not a random-renewal recovery
guarantee. Full population gradients and the deliberately chosen
starting local minimum are part of the experiment.

\section{Independent Gaussian initialization: evidence, not a theorem}
At width 16, 64 independent-Gaussian seeds, and squared population
loss, all runs at scale $10^{-4}$ reached the 50\% uniform-test-error
shortcut by time 2000. At scales $0.01,0.1,0.5$, respectively
54, 12, and 0 of 64 runs had exactly 50\% uniform-test error.
The dynamics use a midpoint Lorentz update that preserves each
$a_j^2-\|w_j\|^2$; the Gaussian runs do not enforce balance.
Step-size and longer-time checks are separate numerical requirements.

A route toward an almost-sure frozen-field result is to write
$a'=U(w)$, $w'=a\nabla U(w)$, whence
$a''=a\|\nabla U(w)\|^2$ almost everywhere. Here the shortcut component
of every frozen gradient has a uniform strictly positive lower bound.
Incoming-dominant neurons may change readout sign. Excluding their
stable set and then transferring their dynamics to the actual finite
network needs a careful nonsmooth-flow argument. It is not claimed here.
\end{document}
